\section{LeanExe: compilation and per-artifact verification}
\label{sec:leanexe}

\subsection{A checked Lean definition as an algorithm specification}

LeanExe imports an elaborated Lean declaration from a built module and compiles its reachable executable definitions.  Lean performs parsing, elaboration, type checking, and proof checking before LeanExe inspects the declaration.  The accepted runtime language is a restricted, pure, monomorphic, first-order subset with machine integers, bounded natural-number computations, packed byte arrays, supported array forms, structures, and internal recursive data.  Supported loops and helper calls express the model's numerical computation.  The language specification identifies accepted primitives and rejects residual expressions outside the subset~\cite{language,subset}.

The source definition supplies a precise specification because it fixes both control flow and arithmetic choices.  In GPT-2, this includes the order of each dot product, the offsets of runtime weight tensors, the byte representation of the cache, and the return value on an invalid input.  A mathematical description such as ``compute attention'' leaves these choices open.  The Lean definition resolves them, and the target proof uses the same definitions when stating expected results~\cite{cached,kernels,numerics}.

Lean's proof language can express richer predicates than the executable subset.  A theorem may quantify over arbitrary initial stores, describe ownership and disjointness, or compare an emitted computation with a mathematical recurrence.  Those proof expressions serve as specifications and evidence.  The compiler emits the runtime-relevant program after erasure of proof arguments and fields with no runtime content.  Its accepted-language checks constrain the computation that survives erasure~\cite{language}.

The theorem subject remains a particular program and specification.  The developer chooses the specification and reviews its meaning.  Lean's kernel checks whether the supplied proof has that type.  These two checks answer different questions: the source definition identifies the intended algorithm, and kernel checking establishes the formal implication.  The statement-review lemmas described later examine whether the public GPT-2 theorem entails concrete successful runs, the complete logit length, and the expected cache growth~\cite{reviewlemmas}.

\subsection{Compilation, representation, and ownership}

The compiler specializes supported source definitions, lowers them into a first-order intermediate representation, and emits structured WebAssembly instructions and their binary encoding.  It also emits WAT, WebAssembly's text format, for inspection.  A library-mode module exports its entry and memory-management functions.  Arrays and byte arrays become objects in linear memory with representations fixed by the ABI.  Scalar instructions and packed loads and stores implement the source primitives~\cite{language,compiler-implementation}.

Packed FP32 storage is central to the GPT-2 implementation.  A 124,439,808-word checkpoint occupies 497,759,232 bytes.  Keeping each value in four bytes permits the full model to fit within a 32-bit WebAssembly memory and provides byte-level views that host code can copy to GPU buffers.  Logical source values retain their \code{ByteArray} interface, while compiler-recognized word operations emit direct reads and construction loops.  The packed-memory proof relates the linear-memory bytes to the source array at every access~\cite{packed,manifest}.

Ownership information determines which allocations a generated operation may retain, return, or release.  The artifact proof checks the resulting memory behavior against representation and heap predicates.  A returned byte array carries both its content and the ownership facts needed by later calls.  The caller can release intermediate buffers only after establishing that the remaining live objects are disjoint from the released allocation.  GPT-2's block and session proofs use those facts to retain weights and current cache while freeing temporary vectors and obsolete cache objects~\cite{entry,budget,session}.

This treatment makes memory management part of functional correctness.  A correct numerical result stored at an invalid or reused pointer would fail the representation postcondition.  A release that overwrites a live tensor would fail the preservation predicate required by the next operation.  Allocation theorems account for either reuse of a suitable free block or growth within the represented memory capacity.  The top-level session derives the resource inequalities used by those alternatives~\cite{budget,input}.

\subsection{Proofs for a compilation}

LeanExe's artifact method proves behavior for an identified compiler output.  The compiler can produce code, annotations, and proposed structural equalities.  Each fact used by the final theorem must follow from checked definitions and proofs.  The current GPT-2 result therefore establishes correctness of its compilation without requiring a theorem quantifying over every possible compiler input~\cite{artifact,binary}.

The binary proof embeds the entire executable byte sequence in Lean.  A proved-sound decoder recovers a formal syntax tree, a proved-sound validator establishes the supported module's validity, and a translation maps that module to Talos.  A checked equality identifies the translated module with the declarations used by the execution proof.  The final theorem substitutes through that equality.  A parsing or translation discrepancy prevents the equality from closing and therefore prevents the artifact theorem from following~\cite{binary,artifact}.

The complete bytes and the file digest have different roles.  The theorem concerns the bytes as a Lean value.  SHA-256 identifies a file in records and helps compare published artifacts.  The proof's semantic conclusion follows from the full byte value, decoding, and execution theorem.  The command that applies the proof to an external file must still establish that the file contains those bytes.  The focused artifact checker performs that comparison~\cite{artifactgate}.

Per-compilation verification also permits implementation changes.  A changed compiler output creates a new proof obligation against the same source algorithm or a revised specification.  Shared arithmetic, memory, and control-flow lemmas can survive the change when their hypotheses remain applicable.  The WGSL development uses this property to replace matrix implementations, preserve proofs for unchanged functions, and reconstruct the controller proof around new call interfaces~\cite{wg-transport,wg-composition}.

\subsection{Annotations and reusable proof structure}

Generated target programs contain repeated instruction patterns for loads, stores, loops, allocation, and stack management.  Compiler annotations identify candidate regions and record their intended source role.  Checked region equalities establish that the instructions at those locations match a proved template.  The proof can then use a theorem about the template while preserving a connection to the emitted program.  The annotation supplies a proposed correspondence that must survive checking~\cite{annotations,proof-method}.

Shared loop rules express a source recurrence and the target state that realizes it.  An invariant may assert that the first $i$ output words are initialized, that an accumulator equals a source fold prefix, and that all protected input bytes retain their initial values.  The rule composes one verified body step through the loop and establishes the final result.  The GPT-2 range-fold rule also preserves an arbitrary operand-stack prefix, which is needed when a destination address remains on the stack during an inner reduction~\cite{loops}.

The repository's broader proof-generation tools use retrieved lemmas, tactics, and worked examples to propose proofs.  Compiler annotations and a structured knowledge library help select applicable material.  The proof-producing process may use coding agents, while the final check accepts Lean proof terms against the exact program and specification.  The GPT-2 proof described here uses the shared proof library and explicit component composition.  Its checked theorem depends on those accepted declarations~\cite{proof-method,ltg,entry}.

This architecture separates the cost of constructing a proof from the basis for accepting it.  An expensive search, a hand-written lemma, or an automatically generated application can produce the same checked proposition.  Reusable results reduce the amount of target-level reasoning needed for each new component.  In this development, scalar arithmetic, packed memory, allocator behavior, loop execution, and function-region transport provide the main reusable interfaces.  Architecture-specific theorems then compose them around fixed GPT-2 dimensions~\cite{f32,packed,budget,wg-transport}.

\subsection{Practical value and scope of the approach}

The result gives a precise answer to an implementation question: whether the emitted module executes the selected algorithm on every input in its specified domain.  It covers mistakes in arithmetic lowering, address calculations, control flow, cache assembly, and memory management when those mistakes would violate the source result or required successful execution.  A checkpoint change of the same byte length preserves the theorem's applicability because the weight values remain quantified inputs~\cite{session,artifact}.

The source remains available for algorithm review.  If a numerical method or cache convention in that source is unsuitable for an application, source agreement preserves that choice.  A separate numerical analysis can compare the source computation with real formulas, and a model evaluation can measure prediction quality.  The present GPT-2 work prioritizes exact implementation agreement.  PyTorch comparisons and text completions supply empirical evidence about the selected model and numerical methods~\cite{tests,testcode}.

The two backends show how the approach extends across execution interfaces.  The CPU path has a complete binary-to-source theorem within the formal WebAssembly semantics.  The shader path has per-compilation statement proofs and exact packed-layout equalities that compose under a host-execution premise.  The unresolved correspondence between physical shader execution and the strict arithmetic model determines the current limit of the GPU guarantee.  The proof records that premise where the external call enters the controller argument~\cite{wg-backend,wg-execution}.
